\label{sec:background}

Incumbency protection has been part of the informal toolkit of legislative redistricting
since at least the Reconstruction era, when state legislatures routinely drew district
lines to protect sitting members from both parties. The practice became more explicit
and more technically refined as mapping technology improved: by the 1990 cycle,
redistricting software enabled mapmakers to geocode incumbent home addresses and
construct districts that guaranteed each sitting member a residential centre within
their new boundaries~\citep{altman2011public}. By the 2010 cycle, the integration of
precinct-level voter registration and election data with incumbent home locations
allowed mapmakers to engineer not just incumbent residence but also the partisan
composition of the district around that residence~\citep{mcghee2014measuring}.

\subsection{Federal Constitutional Doctrine}

The Supreme Court's treatment of incumbency as a redistricting criterion has been
consistently permissive. In \textit{Karcher v.\ Daggett},\footnote{\textit{Karcher
v.\ Daggett}, 462 U.S.\ 725 (1983).} the Court evaluated a New Jersey congressional
map with minor population deviations and held that such deviations were not justified
by the state's asserted interests. Significantly for present purposes, Justice Brennan's
majority opinion catalogued a non-exhaustive list of legitimate state objectives that
might justify minor population deviations, explicitly including ``making districts
compact, respecting municipal boundaries, preserving the cores of prior districts, and
avoiding contests between incumbents.'' This dicta has been consistently read as
endorsing incumbency protection as a permissible redistricting criterion, though the
Court did not hold that it was required~\citep{karcher1983}.

\textit{Burdick v.\ Takushi}\footnote{\textit{Burdick v.\ Takushi}, 504 U.S.\ 428
(1992).} arose in the context of write-in voting restrictions, not redistricting, but
its articulation of the state interest in ``election administration'' stability has
been invoked in redistricting proceedings to justify incumbent protection as promoting
continuity in representation and reducing disruption to constituent-service
relationships~\citep{burdick1992}.

\textit{Cox v.\ Larios}\footnote{\textit{Cox v.\ Larios}, 542 U.S.\ 947 (2004),
affirming 300 F.Supp.2d 1320 (N.D. Ga. 2004).} drew the outer limit: the Court
summarily affirmed a three-judge district court's holding that Georgia's use of
incumbency protection to justify population deviations exceeding $\pm 5\%$ was
unconstitutional. The Court made clear that incumbency is a permissible criterion
that may justify \textit{minor} deviations, not a licence to deviate substantially
from population equality~\citep{cox2004}.

\subsection{The VRA Tension}

Incumbency protection intersects uncomfortably with the Voting Rights Act of 1965 in
states with substantial minority populations. In \textit{Thornburg v.\ Gingles},
\footnote{\textit{Thornburg v.\ Gingles}, 478 U.S.\ 30 (1986).} the Court established
the three-part test for demonstrating vote dilution under VRA \S{}2: the minority group
must be sufficiently large and geographically compact to constitute a majority in a
single-member district, it must be politically cohesive, and the majority must vote
sufficiently as a bloc to defeat minority-preferred candidates~\citep{gingles1986}.

The tension arises when a sitting incumbent from the majority-race community represents
a district with a substantial minority population. Protecting that incumbent's boundary
by preserving the prior district lines may simultaneously preserve a district that the
VRA would require to be redrawn to create a majority-minority district. Conversely,
creating a majority-minority district to satisfy the VRA may require pairing the existing
Anglo incumbent with a minority challenger or eliminating the incumbent's home district
entirely. No federal court has definitively resolved this tension; it is typically
resolved at the remedial stage through negotiation among parties and special masters.

\subsection{Post-\textit{Rucho} State Court Landscape}

Following \textit{Rucho v.\ Common Cause}\footnote{\textit{Rucho v.\ Common Cause},
588 U.S.\ 684 (2019).} and its holding that federal courts lack jurisdiction to adjudicate
partisan gerrymandering claims, incumbent-protection analysis has migrated to state courts.
State constitutional ``free elections'' clauses have provided the primary vehicle: in
\textit{Harper v.\ Hall}\footnote{\textit{Harper v.\ Hall}, 868 S.E.2d 499 (N.C. 2022).}
the North Carolina Supreme Court invalidated the enacted congressional map in part on
the ground that incumbent protection was used as a tool of partisan advantage---protecting
incumbents of one party at the expense of competitive districts for the other party. The
\textsc{bisect} reference distribution provides state courts with exactly the neutral
baseline they need to evaluate whether the incumbency protection in an enacted map is
consistent with geographic necessity or represents deliberate political engineering.
